\section{Analytic-arc series and mismatch controls}\label{app:arcseries}

For the ARC1 control
\(d=(s+s^3,0,-s,0)\) and \(t=c_ts\), exact Gaussian-rational matrix expansion
gives
\begin{align*}
[s^6]K&=\frac{c_t^4}{12}+\frac{ic_t^3}{6},\\
[s^7]K&=\frac{c_t^4}{12}-\frac{ic_t^5}{60}.
\end{align*}
At \(c_t=-2i\), these are \(0\) and \(4/5\).  The calculation uses the full
matrix exponential series through depth seven; it is not obtained by assuming
the claimed survivor.

For the GF1 mixed path \(d=(s,s^3,-s,s^4)\), \(t=c_ts\), the first two tied
raw moment contributions produce
\begin{align*}
[s^6]K&=-\frac{ic_t^3}{6}+\frac{c_t^4}{12},\\
[s^7]K\big|_{c_t=2i}&=-\frac{32}{15}.
\end{align*}
The first expression vanishes at (2i).  These two controls use opposite signs
of the imaginary time scale and are not copies of one another.

The reachable moment--amplitude mismatch can be summarised as follows.
\begin{center}
\begin{tabular}{lll}
\toprule
Parent points & Moment face & Amplitude face\\
\midrule
\((e_1,4),(e_2,3)\) & \(2c_1-c_2\) & \(c_1/12-ic_2/6\)\\
ARC1 tie & \(2c_\e^2c_t^4+\cdots\) & \(c_\e^2c_t^4/12+ic_\e c_t^3/6\)\\
GF1 mixed path & \(-c^3+2c^4\) & \(-ic^3/6+c^4/12\)\\
\bottomrule
\end{tabular}
\end{center}

The first row proves the logical point most cleanly: non-zero scalar conversion
at each exponent preserves support, while unequal \(k\)-dependent conversions
change the face equation.
